%%%%%%%%%%%%%%%%%%%%%%%%%%%%%%%%%%%%%%%%%
% Short Sectioned Assignment
% LaTeX Template
% Version 1.0 (5/5/12)
%
% This template has been downloaded from:
% http://www.LaTeXTemplates.com
%
% Original author:
% Frits Wenneker (http://www.howtotex.com)
%
% License:
% CC BY-NC-SA 3.0 (http://creativecommons.org/licenses/by-nc-sa/3.0/)
%
%%%%%%%%%%%%%%%%%%%%%%%%%%%%%%%%%%%%%%%%%

%----------------------------------------------------------------------------------------
%	PACKAGES AND OTHER DOCUMENT CONFIGURATIONS
%----------------------------------------------------------------------------------------

\documentclass[paper=a4, fontsize=11pt]{scrartcl} % A4 paper and 11pt font size

\usepackage[T1]{fontenc} % Use 8-bit encoding that has 256 glyphs
\usepackage{fourier} % Use the Adobe Utopia font for the document - comment this line to return to the LaTeX default
\usepackage[english]{babel} % English language/hyphenation
\usepackage{amsmath,amsfonts,amsthm} % Math packages
\usepackage{slashed}

\usepackage{graphicx}
\usepackage{multicol}
\usepackage{enumitem}
\usepackage{caption}
\usepackage{esint}

\usepackage{listings}
\usepackage{color}

\definecolor{dkgreen}{rgb}{0,0.6,0}
\definecolor{gray}{rgb}{0.5,0.5,0.5}
\definecolor{mauve}{rgb}{0.58,0,0.82}

\lstset{frame=tb,
  language=C,
  aboveskip=3mm,
  belowskip=3mm,
  showstringspaces=false,
  columns=flexible,
  basicstyle={\small\ttfamily},
  numbers=none,
  numberstyle=\tiny\color{gray},
  keywordstyle=\color{blue},
  commentstyle=\color{dkgreen},
  stringstyle=\color{mauve},
  breaklines=true,
  breakatwhitespace=true
  tabsize=3
}

\usepackage{lipsum} % Used for inserting dummy 'Lorem ipsum' text into the template

\usepackage{sectsty} % Allows customizing section commands
\allsectionsfont{\centering \normalfont\scshape} % Make all sections centered, the default font and small caps

\usepackage{abstract}
\renewcommand{\abstractnamefont}{\normalfont\Large}
\renewcommand{\abstracttextfont}{\normalfont}

\usepackage{fancyhdr} % Custom headers and footers
\pagestyle{fancyplain} % Makes all pages in the document conform to the custom headers and footers
\fancyhead{} % No page header - if you want one, create it in the same way as the footers below
\fancyfoot[L]{} % Empty left footer
\fancyfoot[C]{} % Empty center footer
\fancyfoot[R]{\thepage} % Page numbering for right footer
\renewcommand{\headrulewidth}{0pt} % Remove header underlines
\renewcommand{\footrulewidth}{0pt} % Remove footer underlines
\setlength{\headheight}{13.6pt} % Customize the height of the header

\numberwithin{equation}{section} % Number equations within sections (i.e. 1.1, 1.2, 2.1, 2.2 instead of 1, 2, 3, 4)
\numberwithin{figure}{section} % Number figures within sections (i.e. 1.1, 1.2, 2.1, 2.2 instead of 1, 2, 3, 4)
\numberwithin{table}{section} % Number tables within sections (i.e. 1.1, 1.2, 2.1, 2.2 instead of 1, 2, 3, 4)

%\setlength\parindent{0pt}  Removes all indentation from paragraphs - comment this line for an assignment with lots of text

%Average
\newcommand{\average}[1]{\ensuremath{\left\langle #1 \right\rangle}}
%Parenthesis 
\newcommand{\parentheses}[1]{\ensuremath{\left( #1 \right)}}
%Commutator
\newcommand{\commutator}[1]{\ensuremath{\left[ #1 \right]}}
%Anti-commutator
\newcommand{\anticommutator}[1]{\ensuremath{\left\lbrace #1 \right\rbrace}}
%QED
\newcommand{\QED}{\begin{flushright}\textit{Q.E.D.}\end{flushright}}
%Split
\newcommand{\spliteq}[1]{\begin{split} #1 \end{split}}
%----------------------------------------------------------------------------------------
%	TITLE SECTION
%----------------------------------------------------------------------------------------

\newcommand{\horrule}[1]{\rule{\linewidth}{#1}} % Create horizontal rule command with 1 argument of height

\title{
\vspace{-2.5cm}
\begin{center}
\includegraphics[width=2.5cm]{logo-kth.png}\\[-1mm]
\hspace{-3mm}
\end{center}
\normalfont \normalsize
\textsc{Theoretical Physics} \\ [25pt] % Your university, school and/or department name(s)
\horrule{0.5pt} \\[0.4cm] % Thin top horizontal rule
\huge Quantum Field Theory\\ % Course name
SI2410\\ % Course code
\huge Fall 2014\\ % Semester
\huge Preparation Questions \\ % The assignment title
\horrule{2pt} \\[0.5cm] % Thick bottom horizontal rule
}

\author{The whole group}
\date{\normalsize\today} % Today's date or a custom date

\begin{document}

\maketitle % Print the title

%----------------------------------------------------------------------------------------

\section{Seminar 1: Functional integrals and introduction to renormalization}

\subsection{What is the relationship between the functional integral formalism and
the n-point correlation function?}
General generating functional is given by

\begin{equation}
Z[J]=\int D\phi \exp \left[i\int d^4x[L+J(x)\phi(x)] \right]
\end{equation}

where $L$ is the Lagrangian density and $J(x)\phi(x)$ is a source term.

The n-point correlation function is then given by

\begin{equation}
\left\langle 0 | T \phi (x_1) ... \phi (x_n) | 0 \right\rangle = \dfrac{1}{Z_0} \prod_{i=1}^{n} \parentheses{-i\dfrac{\delta}{\delta J(x_i)} Z[J]}
\end{equation}

where $Z_0=Z[J=0]$. This yields

\begin{equation}
\left\langle 0 | T \phi (x_1) ... \phi (x_n) | 0 \right\rangle = \dfrac{\int D\phi \phi(x_1)...\phi(x_n) \exp \left[i\int d^4xL \right]}{\int D\phi \exp \left[i\int d^4xL \right]}
\end{equation}

\subsection{Given a Lagrangian density, how can the Feynman rules of a theory be
computed using functional integral formalism?}
The Feynman rules can be computed by Taylor expanding

\begin{equation}
\exp iS = \exp \commutator{i \int d^4 x L_0 + L_{int}}
\end{equation}

where $S$ is the action, $L_0$ is the free Lagrangian density and $L_{int}$ is the interaction Lagrangian density. In the case of $\lambda \phi^4$-theory we have the Lagrangian

\begin{equation}
L=L_0-\dfrac{\lambda}{4!}\phi^4
\end{equation}

and by Taylor expansion

\begin{equation}
\exp iS = \exp \commutator{i \int d^4 x L_0 - \dfrac{\lambda}{4!}\phi^4}= \exp \commutator{i \int d^4 x L_0} \parentheses{1-i\int d^4x \dfrac{\lambda}{4!}\\phi^4+...}
\end{equation}

where $L_0$ gives the propagators and all terms of order higher than two yields interaction. In this case we read of the vertex factor in momentum space as

\begin{equation}
-i\lambda (2 \pi)^4 \delta ^{4} \parentheses{\sum p}
\end{equation}

\subsection{What complications arise when using the functional integral formalism to
quantize the electromagnetic field? How is it solved?}

\subsection{What are the properties of Grassmann numbers and how are they used to
quantize spinor fields?}

\subsection{What is the superficial degree of divergence and how can it be computed?
Use QED as an example.}

\subsection{How does renormalized perturbation theory relate the bare and physical
masses?}

\end{document}